\booksection{第十一章付録：プロフェッショナル映像技術知識}

\chapterintro{スマホ撮影と動画の奥義を真に理解するには、いくつかの専門的な映像技術知識を理解する必要があります。これらの知識はスマホカメラをより良く使うのに役立つだけでなく、デバイス選びの際の判断力も高めてくれます。}

\booksubsection{一、解像度、フレームレート、ビットレート：動画品質の三要素}

\subsubsection*{解像度は画面の鮮明度を決定し、ピクセル数で表されます。}

720P (HD)

\begin{itemize}
  \item \textbf{：1280×720ピクセル、エントリーレベルの鮮明度。}1080P (Full HD)
  \item \textbf{：1920×1080ピクセル、現在主流の解像度。}2K (QHD)
  \item \textbf{：2560×1440ピクセル、1080Pより鮮明。}4K (UHD)
  \item \textbf{：3840×2160ピクセル、超高画質、ディテール豊富。}8K
  \item \textbf{：7680×4320ピクセル、未来の標準、現状は対応デバイスが限定的。}解像度が高いほど画面は鮮明ですが、ファイルも大きくなり、ストレージと処理能力の要求が高まります。
\end{itemize}

フレームレート（Frame Rate）：

\subsubsection*{フレームレートは画面の滑らかさを決定し、毎秒フレーム数（fps）で表されます。}

24fps

\begin{itemize}
  \item \textbf{：映画標準フレームレート、独特の「シネマ感」があります。}30fps
  \item \textbf{：テレビやネット動画で一般的、滑らかで自然。}60fps
  \item \textbf{：高フレームレート、スポーツシーンやゲーム録画に適し、非常に滑らか。}120fps/240fps
  \item \textbf{：超高フレームレート、スローモーション動画に使用可能。}フレームレートが高いほど画面は滑らかで、スローモーション効果も良くなりますが、ファイルも大きくなります。
\end{itemize}

ビットレート（Bitrate）：

\subsubsection*{ビットレートは動画のデータ密度を決定し、毎秒データ量（Mbps）で表されます。}

ビットレートが高いほど画質が良く、ディテールが多く保持されますが、ファイルも大きくなります。

\begin{itemize}
  \item ビットレートが低すぎると、画面にモザイクや色塊などの圧縮痕が現れます。
  \item プロフェッショナルな動画撮影では通常、高いビットレートが必要で、後処理の余地を多く残せます。
  \item 二、動画コーデック形式：H.264、H.265とProRes
\end{itemize}

\booksubsection{二、動画コーデック形式：H.264、H.265とProRes}

\subsubsection*{特徴}

\begin{itemize}
  \item \textbf{特徴}長所
  \item \textbf{：後処理の調整余地が極めて大きく、露出やホワイトバランスなどのパラメータを画質劣化なしに変更可能。}短所
  \item \textbf{欠点}適用シーン
  \item \textbf{：ハイエンドプロフェッショナル制作。iPhone 15 ProシリーズはApple ProRes動画撮影に対応（最大4K 60fps）しますが、ProRes RAW形式は現状外部録画デバイスとの組み合わせが必要で、スマホ単体での直接撮影は未対応。}H.265 (HEVC)：
\end{itemize}

\subsubsection*{特徴}

\begin{itemize}
  \item \textbf{特徴}長所
  \item \textbf{：後処理の調整余地が極めて大きく、露出やホワイトバランスなどのパラメータを画質劣化なしに変更可能。}短所
  \item \textbf{欠点}適用シーン
  \item \textbf{：ハイエンドプロフェッショナル制作。iPhone 15 ProシリーズはApple ProRes動画撮影に対応（最大4K 60fps）しますが、ProRes RAW形式は現状外部録画デバイスとの組み合わせが必要で、スマホ単体での直接撮影は未対応。}Apple ProRes：
\end{itemize}

\subsubsection*{特徴}

\begin{itemize}
  \item \textbf{特徴}長所
  \item \textbf{：後処理の調整余地が極めて大きく、露出やホワイトバランスなどのパラメータを画質劣化なしに変更可能。}短所
  \item \textbf{欠点}適用シーン
  \item \textbf{：ハイエンドプロフェッショナル制作。iPhone 15 ProシリーズはApple ProRes動画撮影に対応（最大4K 60fps）しますが、ProRes RAW形式は現状外部録画デバイスとの組み合わせが必要で、スマホ単体での直接撮影は未対応。}ProRes RAW：
\end{itemize}

\subsubsection*{特徴}

\begin{itemize}
  \item \textbf{特徴}RAWとは
  \item \textbf{：未処理の生センサーデータ、最大情報量を保持。}長所
  \item \textbf{：後処理の調整余地が極めて大きく、露出やホワイトバランスなどのパラメータを画質劣化なしに変更可能。}短所
  \item \textbf{欠点}適用シーン
  \item \textbf{：ハイエンドプロフェッショナル制作。iPhone 15 ProシリーズはApple ProRes動画撮影に対応（最大4K 60fps）しますが、ProRes RAW形式は現状外部録画デバイスとの組み合わせが必要で、スマホ単体での直接撮影は未対応。}：ハイエンド専門制作向け。iPhone 15 Proシリーズ以降はApple ProRes動画撮影に対応（最大4K 60fps）、iPhone 17 ProはさらにProRes RAWの本体直接録画に対応しています。
\end{itemize}

\booksubsection{三、HDR：ハイダイナミックレンジ技術}

HDR10

\begin{itemize}
  \item \textbf{：最も基本的なHDR標準、互換性が良い。}HDR10+
  \item \textbf{：Samsungが推進する標準、ダイナミックメタデータに対応。}Dolby Vision
  \item \textbf{：Dolby開発のトップHDR標準、画質が最も良いがライセンスが必要。}HLG
  \item \textbf{：放送級HDR標準、通常モニターとの互換性あり。}HDR動画撮影時は、HDR対応ディスプレイが必要で完全な効果を確認できます。
\end{itemize}

四、色空間と色深度

\booksubsection{四、色空間と色深度}

\subsubsection*{sRGB}

\begin{itemize}
  \item \textbf{：標準色空間、互換性が最も良い。}Display P3
  \item \textbf{：Appleが推進する広色域標準、色がより豊か。}Rec.2020
  \item \textbf{：HDR動画の標準色空間、色域が最も広い。}色深度：
\end{itemize}

\subsubsection*{8-bit}

\begin{itemize}
  \item \textbf{：1670万色、通常動画標準。}10-bit
  \item \textbf{：10.7億色、プロフェッショナル動画標準、色のグラデーションがより滑らか。}12-bit
  \item \textbf{：680億色、ハイエンドプロフェッショナル標準。}実用的なアドバイス
\end{itemize}

\begin{tcbinfo}{日常撮影：1080P/4K + 30fps + H.265エンコード、画質とファイルサイズのバランス。}
\begin{itemize}
  \item スポーツ撮影：1080P/4K + 60fps、滑らかさを確保。
  \item プロフェッショナル創作：4K + ProRes（対応場合）、最大の後処理余地を保持。
  \item ストレージの提案：高ビットレート動画は容量を食うので、定期的にパソコンやクラウドストレージにバックアップ。
  \item 探索してみよう：スマホの動画設定
\end{itemize}

\end{tcbinfo}

\begin{tcbpractice}{カメラAppを開き、動画モードに切り替え。}
\begin{enumerate}
  \item 解像度とフレームレートのオプションを確認し、異なる組み合わせを試す。
  \item 対応していれば、HDRやProResなどの高度なオプションがあるか確認。
  \item 異なる設定で数本の動画を撮影し、画質とファイルサイズを比較。
  \item 第十二章：映像エンターテインメントと読書
\end{enumerate}

\end{tcbpractice}

\booksubsection{五、RAWフォーマットとHEIF/JPEGの比較}

\subsubsection*{RAWフォーマットとは：}

RAWフォーマットはカメラセンサーが記録した「生データ」——一切の圧縮や処理を経ていない純粋な原始情報です。生鮮食品のように、後処理で「調理」して完成させます。スマホはストレージと処理能力が限られているため、デフォルトで撮影されるのはすでに「調理済み」のJPEGまたはHEIFで、RAWは特地に有効にする必要があります。

\subsubsection*{RAW vs HEIF vs JPEGの比較：}

\begin{itemize}
  \item \textbf{JPEG}：最も古いフォーマットで、互換性が最も高い（どのデバイスでも開ける）。8-bit色深度、非可逆圧縮、ファイルサイズが最小（1枚2-5MB）。後処理の余地が限られており、無理に暗部を明るくすると色ノイズが出やすいです。
  \item \textbf{HEIF/HEIC}：Appleが推進する新フォーマットで、同じ画質でファイルサイズはJPEGの半分（1枚1-3MB）。10-bit色深度と透明チャネルに対応。欠点は一部の旧デバイスやWindowsではプラグインのインストールが必要なことです。
  \item \textbf{RAW}：非圧縮の生データで、ファイルサイズが最大（1枚20-100MB）。色深度は通常12-16bitで、後処理の余地が最大。欠点は専門ソフト（Lightroom、Snapseedなど）でないと開いて処理できず、直接共有できないことです。
\end{itemize}

\subsubsection*{なぜプロカメラマンはRAWで撮るのか：}

\begin{itemize}
  \item \textbf{ホワイトバランスを後処理で無劣化調整できる}：JPEGでホワイトバランスを間違えると（例：白熱灯下で黄色く偏る）、後処理で補正すると画質が劣化します。RAWは色温度を再設定でき、完全に無劣化です。
  \item \textbf{露出許容範囲が広い}：JPEGで2段アンダー露出の暗部はほぼ真っ黒で、明るくするとノイズだらけです。RAWは3-4段アンダーの暗部から「救える」使用可能なディテールを取り出せます。過曝も同様です。
  \item \textbf{暗部とハイライトのディテールがより豊か}：RAWのダイナミックレンジはJPEGよりはるかに大きく、明部が飛ばず、暗部が潰れない余裕がずっと大きいです。
  \item \textbf{シャープ化、ノイズ除去、レンズ補正がすべて後処理で制御可能}：カメラの本体処理は「一発勝負」ですが、RAWは何度でも調整でき、毎回生データから出発できます。
\end{itemize}

\subsubsection*{スマホでRAWを有効にする方法：}

\begin{itemize}
  \item \textbf{iPhone ProRAW}：iPhone 12 Pro以上のPro機種のみ対応。「設定 カメラ フォーマット」で「Apple ProRAW」をオンにします。カメラ画面上部に「RAW」ボタンが表示されるのでタップして有効化。ファイルサイズは約25MB/枚、必要な時のみオンにすることをお勧めします。
  \item \textbf{Android Motion Raw / Ultra RAW}：ブランドによって名称が異なります。Xiaomi/OPPO/vivoのフラッグシップ機はカメラの「プロモード」または「その他」でRAWスイッチが見つかります。Huaweiは「プロモード」でRAWを有効化。Samsung Galaxy Sシリーズは「プロ」モードにRAWオプションがあります。ファイルサイズは通常20-50MB。
  \item \textbf{一般的なアドバイス}：日常のスナップ撮影でRAWを有効にする必要はありません。ファイルが大きく容量を占めます。重要なシーン（風景、ポートレート、後処理が必要な場合）のみ有効にしましょう。
\end{itemize}

\booksubsection{六、コンピュテーショナルフォトグラフィの原理}

現代のスマホ撮影は「コンピュテーショナルフォトグラフィ」なしでは成り立ちません——つまり、アルゴリズムによってレンズやセンサーだけでなく写真を生成します。これらの核心原理を理解すると、スマホがシャッターを切った瞬間に何をしているかがよりわかります。

\subsubsection*{マルチフレーム合成：}

スマホでシャッターを一度切ると、実は複数枚の写真を撮っています。アルゴリズムが複数フレームを整列させ、加重融合して高品質な一枚を作ります。

\begin{itemize}
  \item \textbf{ナイトモード}：1-3秒で複数枚のアンダー露出写真を連写し、整列後に重ね合わせ——ノイズが互いに相殺され、明るさが累積します。結果として暗所でもクリーンな画像が撮れます。
  \item \textbf{HDR（ハイダイナミックレンジ）}：アンダー、正常、オーバーの3枚（または複数枚）を撮影し、明部と暗部の両方にディテールがある1枚に合成します。逆光シーンで特に効果的です。
  \item \textbf{Deep Fusion（Apple）/Photonic Engine（光子エンジン）}：中低光シーン向けに、ピクセル単位で複数フレームを分析·融合し、テクスチャとディテールを重点的に向上させます。セーターや髪の毛などのディテールが明らかに鮮明になります。
\end{itemize}

\subsubsection*{セマンティックセグメンテーション：}

AIは単に画像全体を一律に処理するのではなく、まず画面に何があるかを「理解」します——空、顔、芝生、建物、衣服——そして各領域に異なる処理戦略を適用します。

\begin{itemize}
  \item \textbf{空}：青色を個別に強調し、明度を下げて過曝を防ぎます。
  \item \textbf{顔}：肌色領域を認識し、個別に明るくし、肌を滑らかにします。目鼻立ちや髪のディテールはいじりません。
  \item \textbf{芝生/葉}：緑の彩度を強調し、ディテールのシャープネスを向上させます。
  \item \textbf{背景ぼかし}：ポートレートモードで、AIが被写体の輪郭を認識し、背景をぼかして被写体を鮮明にします。
\end{itemize}

\subsubsection*{ディープラーニングの応用：}

\begin{itemize}
  \item \textbf{AIシーン認識}：カメラが「グルメ」「ペット」「夜景」「書類」などのシーンを自動認識し、対応するパラメータを最適化します。例えばグルメモードは暖色と彩度を強調し、料理がより魅力的に見えます。
  \item \textbf{AIポートレートぼかし}：スマホレンズは物理的に一眼レフのような浅い被写界深度を実現できないため、ディープラーニングモデルで被写体の輪郭を予測し、背景をぼかします。最新アルゴリズムは髪の毛やガラスの端などのディテールも正しく処理できます。
  \item \textbf{AI消去/AI拡張}：画面内の通行人や不要物を選択すると、AIがコンテキストに基づいて隠された部分を「補完」します。これは生成AIの映像分野での応用です。
\end{itemize}

\subsubsection*{コンピュテーショナルフォトグラフィ vs 光学撮影：}

\begin{itemize}
  \item \textbf{コンピュテーショナルフォトグラフィの優位性}：小型でも高画質が得られ、暗所や逆光などの過酷なシーンでも撮影でき、一般ユーザーが写真の知識がなくても「綺麗な」写真が撮れます。
  \item \textbf{コンピュテーショナルフォトグラフィの弱点}：画像が時に「過度処理」になる——過剰シャープ化、色の歪み、不自然な空の色。被写体が動くとマルチフレーム合成でゴーストが残る。ディテールは「計算で作られた」ものであり、真に記録されたものではない。
  \item \textbf{光学撮影の優位性}：大型センサー+大口径レンズで光をリアルに捉え、画質が自然で後処理の余地が大きい。アルゴリズムの「補完」によるアーティファクトが発生しない。
  \item \textbf{トレンド}：現代のフラッグシップスマホは「光学+コンピュテーショナル」の融合——大型センサーでリアルな情報を取得し、アルゴリズムで表示を最適化。両者は対立するものではありません。
\end{itemize}

\booksubsection{七、レンズタイプと焦点距離}

現代のフラッグシップスマホは通常3〜4つのレンズを搭載し、それぞれが異なる焦点距離に対応します。焦点距離を理解すると、シーンに応じたレンズ選択ができます。

\subsubsection*{一般的なレンズと換算焦点距離：}

\begin{itemize}
  \item \textbf{超広角（13-24mm換算）}：視野が最も広く、風景、建築、屋内空間に適しています。目の前の大シーンを撮影できますが、エッジに歪みがあり（直線が曲がる）、ポートレートでは顔が引き伸ばされます。「大シーン」に適し、「クローズアップ」には不向きです。
  \item \textbf{メインカメラ（24-28mm換算）}：日常で最もよく使うレンズで、画質が最も良く、絞りが最大。視野と画質を両立し、人物も風景も撮れます。ほとんどのスマホの「1x」がメインカメラです。
  \item \textbf{望遠（70-200mm換算）}：遠景やポートレートのクローズアップに使用。望遠は空間の圧縮効果があり（遠くの山と近くの人がより近く見える）、ポートレートではより自然な目鼻立ちの比率と背景ぼかしが得られます。
  \item \textbf{ペリスコープ式望遠（200mm+換算）}：反射式光路設計により、スマホの厚み内でより長い焦点距離を実現。遠距離でもディテールをクリアに撮影でき、コンサート、スポーツイベント、野生動物に適しています。フラッグシップ機のみ搭載。
\end{itemize}

\subsubsection*{センサーサイズと画素サイズのトレードオフ：}

センサーは「ベース」であり、受け取れる光量を決定します。業界では「ベースが一段大きいと圧倒的」と言われます——同じ画素数なら、センサーが大きいほど1画素あたりのサイズも大きく、集光量が多く、画質が良い。

\begin{itemize}
  \item \textbf{大型センサー}：集光量が多く、夜景の画質が良く、ダイナミックレンジが広い。欠点はスマホが厚くなり（カメラモジュールが突出）、コストが高いこと。
  \item \textbf{高画素·小型センサー}：昼間の光が十分な時は、高画素でより鮮明な写真が撮れ、後処理の切り抜きに便利。ただし暗所では1画素あたりの集光が不足し、ノイズが多くなる。
  \item \textbf{主流フラッグシップの選択}: メインカメラは1/1.3インチから1インチ（一部機種は1インチ大型センサーを採用）、画素数は約5000万——これが画質と体積のバランスポイントです。
\end{itemize}

\subsubsection*{画素融合技術（Pixel Binning）：}

昼間の高画素と夜間の大画素を両立するため、スマホメーカーは「画素融合」を使用——複数の小画素を1つの大画素に統合して出力します。

\begin{itemize}
  \item \textbf{4in1画素（Quad Bayer）}：5000万画素センサーで1250万画素写真を出力。各出力画素は4つの隣接画素から合成され、集光量が4倍向上し、画質がよりクリーンに。
  \item \textbf{9in1画素（Nonacell）}：1億画素以上のセンサーに使用し、約1100-1200万画素を出力。同じ原理で、集光量がさらに大きい。
  \item \textbf{実際の効果}：デフォルト撮影モードでは、融合後の低画素写真が得られます——画質が良く、ファイルが小さい。究極のディテールが必要な時（書類撮影や遠景など）は、手動で「高画素モード」に切り替えて全解像度出力できます。
\end{itemize}

\booksubsection{八、特殊撮影モード}

日常撮影以外に、スマホにはいくつか特殊モードがあり、使いこなせば普段は撮れない画面が撮れます。

\subsubsection*{パノラマモード：}

\begin{itemize}
  \item 撮影時に案内矢印に従ってスマホをゆっくり水平に移動してください。速すぎると——継ぎ目がずれます。
  \item スマホの垂直方向を安定させ、水平にのみ移動してください。上下に揺れると水平線が波状になります。
  \item 動く人や車がある場合は注意——合成時に「体が半分」や「ゴースト」が現れます。
  \item パノラマモードは縦撮影も可能——縦に撮って横に合成し、高い建物や滝に適しています。
\end{itemize}

\subsubsection*{スローモーションとタイムラプス：}

\begin{itemize}
  \item \textbf{スローモーション}：120fpsまたは240fpsの高フレームレートで撮影し、再生時にスロー再生。水滴、スポーツの瞬間、ペットのジャンプに適しています。光が十分必要で、スローモーションモードは露出時間が短く、暗所では画質が低下します。
  \item \textbf{タイムラプス}：間隔撮影で複数枚の写真を合成し早送り動画を作成。車の流れ、雲の動き、日落、花の開花に適しています。撮影時間は数分から数時間で、必ず三脚でスマホを固定してください。手持ちでは確実にブレます。
  \item \textbf{テクニック}：タイムラプスで雲を撮る場合、スマホを固定して15-30分撮影すると、最終的に約10秒のタイムラプス動画が得られ、インパクトがあります。
\end{itemize}

\subsubsection*{星空·天文モード：}

\begin{itemize}
  \item 都市の光害から遠ざかり、街灯の干渉がない郊外や山間部を見つけてください。都市では天の川は撮れません。
  \item 必ず三脚を使用し、手持ちは不可。スマホを数秒から数十秒しっかり固定する必要があります。
  \item 天気は晴れて雲がないこと、月が大きすぎないこと（満月は星光を隠します）。月初めや月末の暗い夜が最適です。
  \item iPhoneの「ナイトモード」は星空を自動認識して露出時間を延長（最大30秒）。Androidフラッグシップの「星空モード」や「天文モード」も同様の効果。辛抱強く待ち、シャッターを切った後はスマホを動かさないでください。
\end{itemize}

\subsubsection*{マクロモード：}

\begin{itemize}
  \item 花、昆虫、宝石、文字のクローズアップなど小さな被写体の撮影に使用。1-2cmの距離まで寄ってピントが合わせられます。
  \item 一部のスマホには専用マクロレンズ（多くは200万画素または500万画素の小型センサー）があり、画質は一般的。フラッグシップ機は超広角レンズでマクロを兼用し、画質がより良い。
  \item 寄って撮影する際はレンズが被写体に非常に近く、光を遮るため、必要に応じてデスクライトや懐中電灯で補光してください。
  \item ピント距離が極めて近く、被写界深度が非常に浅いため、被写体が前後に数ミリ動くだけでピントがずれる可能性があります。撮影時は手を安定させ、被写体を動かさないのがベスト。
\end{itemize}

\begin{tcbinfo}{静止画撮影の実用アドバイス}
\begin{itemize}
  \item 日常撮影：HEIF/JPEGのデフォルトで十分。ファイルが小さく、画質も実用十分。
  \item 重要なシーン（風景、ポートレート、後処理が必要）：RAWを有効にし、最大の後処理余地を確保。
  \item レンズ選び：風景は超広角、ポートレートは望遠、日常はメインカメラ。
  \item 夜景と星空：必ずスマホを固定し、ナイトモードまたは天文モードを有効化。
  \item 大型センサー搭載スマホ（1インチセンサー機など）は画質の上限が高いですが、操作習慣がやはり鍵。
\end{itemize}

\end{tcbinfo}

\begin{tcbpractice}{試してみよう：あなたのカメラの能力を把握する}
\begin{enumerate}
  \item スマホがRAW撮影に対応しているか確認（設定カメラフォーマット）。有効にして1枚撮影し、同じシーンのJPEGとRAWのファイルサイズを比較。
  \item 超広角、メイン、望遠の3つの焦点距離に切り替え、それぞれ同じ遠景を撮影して視角と画質の違いを体験。
  \item ナイトモードまたは星空モードがあるシーンを見つけ、スマホを固定して夜景写真を1枚撮影。
  \item パノラマモードで横方向に大シーンを撮影してみる。スマホを安定して移動させることに注意。
\end{enumerate}

\end{tcbpractice}
